\section{Grilla local extendida por odometría}

\subsection{Idea general}

El nodo propuesto implementa una \textbf{grilla local extendida egocéntrica}. A diferencia del
\texttt{local\_mapper}, no busca construir un mapa de ocupación probabilístico ni mantener
una representación métrica persistente del entorno. Su objetivo es más acotado: extender
la apertura angular útil de la grilla de profundidad instantánea usando una memoria temporal
corta de mediciones anteriores, reubicadas con la odometría reciente.

La cámara entrega una medición real dentro de su campo de visión frontal, aproximadamente
entre \(-70^\circ\) y \(+70^\circ\). La grilla extendida publica una salida angular de
aproximadamente \(270^\circ\), cubriendo \(-135^\circ\) a \(+135^\circ\). Los sectores
centrales se completan con la medición actual de profundidad; los sectores laterales
\([-135^\circ,-70^\circ]\) y \([70^\circ,135^\circ]\) se completan con mediciones previas,
transformadas al punto de vista actual mediante la odometría.

En palabras simples: el sistema recuerda durante pocos segundos lo que la cámara acaba de
ver, y lo desplaza/rota según el movimiento estimado del usuario para que siga apareciendo
en la grilla local aunque haya salido del campo visual instantáneo.

\subsection{Problema que se busca resolver}

La grilla de profundidad directa tiene una limitación geométrica: sólo representa lo que
está delante de la cámara. En navegación corporal o asistencia háptica esto produce una
zona ciega lateral. Un obstáculo detectado al frente puede quedar fuera de la grilla apenas
el usuario rota la cabeza o avanza, aunque siga siendo relevante para la decisión inmediata.

El problema no requiere necesariamente un mapa global. Para una interfaz háptica reactiva,
esta memoria no debe reemplazar la señal háptica normal: la grilla frontal que recibe el
usuario debe seguir proviniendo de la medición actual. La grilla extendida funciona como
señal auxiliar para advertir obstáculos laterales cercanos fuera del FOV instantáneo. Alcanza
con preservar información local durante una ventana corta, siempre que:

\begin{itemize}
  \item la odometría sea suficientemente coherente en el corto plazo;
  \item las mediciones viejas expiren rápido;
  \item los saltos de odometría invaliden la memoria lateral;
  \item la salida conserve el formato de grilla que ya consume el resto del sistema.
\end{itemize}

Por eso este nodo apunta a robustez local, no a precisión milimétrica. El criterio de éxito
no es reconstruir exactamente el entorno, sino evitar que obstáculos recientemente medidos
desaparezcan de la memoria lateral por estar fuera del FOV actual.

\subsection{Relación con enfoques existentes}

La idea se parece parcialmente a tres familias de métodos:

\begin{itemize}
  \item \textbf{Costmaps locales con ventana móvil}: en Nav2, el \textit{local costmap}
  puede operar como una ventana local en \texttt{odom}, actualizada con observaciones de
  sensores. La documentación de Nav2 describe el uso de capas de obstáculos y ventanas
  locales para navegación de corto plazo. Ver:
  \url{https://github.com/ros-navigation/navigation2/blob/main/nav2_costmap_2d/README.md}
  y \url{https://docs.nav2.org/setup_guides/sensors/mapping_localization.html}.

  \item \textbf{Fusión temporal de profundidad}: existen trabajos donde mapas de
  profundidad consecutivos se fusionan usando transformaciones obtenidas por odometría
  visual, incluso desde un punto de vista egocéntrico, para mejorar percepción reactiva
  sin construir un mapa global pesado. Un ejemplo es ``Gaussian Mixture Models for
  Temporal Depth Fusion'', NASA/JPL:
  \url{https://ntrs.nasa.gov/citations/20210007922}.

  \item \textbf{Mapas egocéntricos}: en robótica móvil se usan representaciones locales
  centradas en el robot para mantener bajo el costo computacional y limitar el efecto de
  deriva. La propuesta de esta tesis toma esa lógica, pero la reduce a una grilla polar
  de profundidad mínima, orientada a salida háptica.
\end{itemize}

La diferencia principal es que aquí no se pretende producir un \texttt{OccupancyGrid} ni
un costmap navegable estándar. Se conserva el contrato de \texttt{DepthGrid}: cada celda
representa una distancia de obstáculo, no una probabilidad de ocupación ni un costo de
planificación.

\subsection{Factibilidad}

La propuesta es factible si se cumplen condiciones moderadas sobre la odometría. No hace
falta una pose precisa a largo plazo, porque la memoria de la grilla es corta. Sí hace falta
que, durante uno o dos segundos, la orientación en yaw y la traslación horizontal sean
coherentes. El error admisible depende de la resolución angular de la grilla: con una salida
de 18 columnas en \(270^\circ\), cada columna cubre aproximadamente \(15^\circ\). Un error
de yaw de pocos grados todavía cae dentro de la misma celda o de una celda vecina; un error
de decenas de grados vuelve insegura la memoria lateral.

El pipeline actual de odometría ya tiene elementos adecuados para este objetivo:

\begin{itemize}
  \item orientación por IMU con filtro complementario;
  \item corrección visual opcional de yaw mediante RGBD;
  \item traslación visual con PnP y RANSAC;
  \item rechazo de saltos por velocidad y traslación máxima plausible.
\end{itemize}

Para esta aplicación la odometría debe priorizar \textbf{consistencia y recuperación ante
fallos}. Si el tracker visual falla o la pose salta, es preferible borrar la memoria lateral
y publicar sólo la grilla actual antes que conservar obstáculos re-proyectados en posiciones
incorrectas. El nodo implementado incorpora compuertas de edad de odometría, salto máximo de
posición y salto máximo de yaw.

\subsection{Comparación con \texttt{local\_mapper}}

No es estrictamente ``mejor'' que \texttt{local\_mapper}; resuelve otro problema.

\begin{itemize}
  \item \texttt{local\_mapper} produce un mapa de ocupación 2D en \texttt{odom}, con
  log-odds, ray casting y olvido espacial. Es más general y más apto para visualización o
  planificación local.

  \item La grilla extendida produce una \texttt{DepthGrid} angular y liviana para alerta
  lateral auxiliar. No estima espacio libre con ray casting ni infiere ocupación
  probabilística.
\end{itemize}

Para una salida háptica reactiva, la grilla extendida puede complementar a la grilla normal
porque conserva la semántica de distancia por celda y tiene bajo costo computacional. No debe
usarse como reemplazo directo de la señal enviada al hardware en condiciones normales. Para
navegación autónoma o planificación, el mapa de ocupación sigue siendo más apropiado.

\subsection{Diseño implementado}

El nodo \texttt{extended\_depth\_grid} se agregó al paquete \texttt{depth\_grid\_encoder}.
Suscribe:

\begin{itemize}
  \item \texttt{/perception/depth\_grid}: grilla actual de profundidad;
  \item \texttt{/nav\_odom}: odometría en el marco \texttt{odom}.
\end{itemize}

Publica:

\begin{itemize}
  \item \texttt{/perception/extended\_depth\_grid}: grilla extendida de profundidad,
  pensada como señal auxiliar de alerta lateral y no como entrada directa del actuador
  háptico principal.
\end{itemize}

Por cada celda válida de la grilla actual se calcula un punto 2D en el plano horizontal,
usando la distancia mínima de la celda y el ángulo de la columna. Ese punto se guarda en
coordenadas \texttt{odom} junto con su fila, estadística de profundidad y timestamp. Al
componer la salida, cada punto histórico se proyecta desde la pose actual: si cae fuera del
FOV real de la cámara pero dentro del FOV extendido, se usa para completar la columna lateral
correspondiente.

Los parámetros principales son:

\begin{itemize}
  \item \texttt{input\_fov\_deg}: FOV real usado por la grilla de entrada;
  \item \texttt{output\_fov\_deg}: FOV de salida, por defecto \(270^\circ\);
  \item \texttt{history\_ttl\_s}: vida máxima de mediciones históricas;
  \item \texttt{max\_odom\_age\_s}: edad máxima aceptada de la odometría;
  \item \texttt{max\_odom\_jump\_m} y \texttt{max\_odom\_jump\_yaw\_deg}: compuertas
  para descartar memoria ante saltos.
\end{itemize}

\subsection{Riesgos y validación necesaria}

Los riesgos principales son:

\begin{itemize}
  \item \textbf{deriva de yaw}: desplaza obstáculos laterales a columnas incorrectas;
  \item \textbf{objetos dinámicos}: una persona u obstáculo móvil puede quedar recordado
  durante el TTL aunque ya no esté allí;
  \item \textbf{oclusiones}: un obstáculo lateral recordado puede haber sido superado o
  tapado por otro objeto;
  \item \textbf{forma de salida}: si se publican 18 columnas para \(270^\circ\), la salida
  no debe enviarse sin adaptación al hardware existente de 10 columnas. En esta etapa se
  consume desde el visualizador como señal de alerta.
\end{itemize}

La validación recomendada es progresiva:

\begin{enumerate}
  \item probar con bag o cámara fija que el nodo publica grillas con dimensiones correctas;
  \item visualizar \texttt{/perception/extended\_depth\_grid} y confirmar que el centro
  coincide con la grilla original;
  \item rotar lentamente la cámara frente a un obstáculo fijo y verificar que el obstáculo
  migra hacia columnas laterales durante el TTL;
  \item provocar saltos o pérdida de odometría y confirmar que la memoria se limpia;
  \item activar la alerta en el visualizador y verificar que la grilla frontal normal no
  cambia;
  \item evaluar la política de alerta lateral para obstáculos fuera del FOV a menos de
  \(0.5\,\mathrm{m}\), antes de cualquier decisión de integración háptica.
\end{enumerate}
